\section{Global Inference 与 Provider Portfolio}

Cursor 同时运行自研模型和调用 frontier providers。Autocomplete 对 latency 极敏感，需要全球 GPU replica 与 overload behavior；复杂 agent request 更看质量、context 与 tool use，也受 provider rate limit 和版本变化影响。本节将 self-hosted serving 与 multi-provider routing 放在同一个 capacity/evaluation 控制面中。

\subsection{Self-hosted model 的 region、queue 与 fallback}

本节先回答 autocomplete 为何不能只部署一个集中集群。读图时沿 user region、router、GPU region 和 fallback 观察：距离、queue depth、model replica 和 KV cache 都会影响 p99；当本地 region 饱和时，系统要在远端 latency、smaller model 和 graceful degradation 之间选择。

\lecturefigure{11-global-inference.png}{Self-hosted inference 需要按 region 路由、容量与降级设计交互尾延迟。}{本地字幕 03:20--04:20、38:10--41:40；概念重绘。}

\begin{knowledgebox}{读图：Capacity headroom 不是浪费}
Interactive serving 不能把平均 utilization 推到 100\%。Queueing 在饱和附近非线性增长；需要 replica headroom、admission control、deadline 和 preemption。低优先级 batch/re-embedding 可以填谷，但必须在在线流量上升时让出资源。
\end{knowledgebox}

\subsection{Frontier providers：Rate limit、quality、cost 与 failover}

\term{rate limit} 是 provider 对 request/token 并发或时间窗口的限制；\term{provider routing} 按任务、质量、健康、quota、成本和 policy 选择模型。简单 round-robin 会忽略不同模型能力与 context/tool interface，真正可用的 portfolio 需要 normalized request、versioned eval 和 provider-specific fallback。

\lecturefigure{12-provider-routing.png}{Multi-provider inference 需要 task contract、provider state、router 与持续证据。}{本地字幕 38:10--41:40；CursorBench 与后续 Cursor Router 资料。}

\begin{knowledgebox}{读图：Offline eval 与 online outcome 各自缺什么}
Offline benchmark 可重复、适合回归，却不完全反映真实 repo、latency 和 user acceptance；online metric 接近产品结果，却受 selection、UI 和 traffic mix 影响。Router 应同时使用二者，并把 model/version 记录到每次 request 和 incident。
\end{knowledgebox}

\begin{importantbox}{Routing 是约束优化}
对 task $x$ 和候选 model/provider $m$，可抽象为
\[
m^*=\arg\max_m Q(m,x)-\alpha C(m)-\beta L(m)
\]
并满足 privacy、availability、context、tool 和 quota 约束。权重 $\alpha,\beta$ 会随 plan、tenant 和交互路径变化，因此不应只维护一个全局“最佳模型”。
\end{importantbox}

\teachervoice{课堂中的现实感在于，快速增长的客户会持续打电话请求更多 capacity，同时准备多 provider 和自有 GPU。Resilience 不是在供应商失败后临时接第二家，而是平时就保持接口、eval、quota 和数据政策可切换。}

\subsection{本章小结}

Global inference 同时是 queueing、capacity、evaluation 和 supply-chain 问题。Self-hosted 与 frontier provider 不是二选一，而是具有不同控制权、成本和故障模式的组合。

